\newpage
\chapter{Техники за детекција на колизии}
Симулатор на меки тела во реално време е, суштински, циклус кој  се движи низ честичките и одговара на прашањето „дали оваа честичка е во контакт со било што?“ за илјадници честички во секој потчекор. Наивниот одговор е да се тестира контакт  на честичката со секој триаголник на секое тело со кое е можно да има контакт, плус сите други честички од телото за да се провери само-судир. Ова има комплексност од $O(N \cdot (T + N))$ во секој потчекор, каде N е бројот на честички и T е вкупниот број на триаголници на статичните објекти. Ако N е во опсегот на неколку илјади а T во опсегот на десетици илјади, многу брзо го надминуваме буџетот за пресметка во една слика за да постигнеме 60 слики во секунда. 

Стандардното решение е да се подели детекцијата на судири во широка фаза и тесна фаза. Во широката фаза се користат едноставни граници на објектите (како axis-aligned bounding boxes - AABBs) за да се одбие поголемиот дел од паровите што немаат шанса да бидат во контакт. Потоа во тесната фаза се користат попрецизни и поскапи методи за да се тестираат само можните кандидати за колизии.

\section{Забрзување во широка фаза (broad-phase acceleration)}

Постојат четири главни структури за забрзување кои се користат во симулациите во реално време.


\textbf{Sweep-and-prune}\cite{sweep}\cite{cohenSweep} . Се одржуваат три подредени листи со крајни точки на AABB, по една за секоја оска. Две AABB се преклопуваат ако и само ако се преклопуваат на сите три оски. Вметнувањата и бришењата се локални; кај умерено динамични сцени, подредените листи ретко се менуваат (претпоставка за временска кохерентност). SAP се истакнува кај цврсти тела кои се движат умерено помеѓу сликите; плаќа трошок за прераспоредување кога геометријата агресивно се деформира.

\textbf{Хиерархиски структури на ограничувачки волумени (BVH)}. Бинарно дрво од вградени AABB (или OBB / сфери / k-DOP)  каде што секој лист опфаќа по еден примитив. Пребарувањето за најблизок објект до некоја  е спуштање со DFS кое ги отстранува сите поддрва чија AABB е веќе подалеку од тековната најдобра вредност. Трошокот за градење е $O(T log T)$, а просечниот трошок за пребарување е $O(log T)$. Oвој пристап е природниот избор за статична геометрија која се пребарува многупати.
	
\textbf{Униформно просторно хаширање }\cite{teschner2003optimized}. 3-Д просторот се мапира на регуларна мрежа од ќелии со страна c, се хешира целобројната координатата на ќелијата во табела со фиксна големина и секој примитив се вметнува во своја кофа на ќелијата. Пребарувањата по честичка ги испитуваат 27-те околни ќелии ($3^3$ за да се обработат граничните ефекти). Вметнувањето е $O(1)$; барањето е $O(k)$, каде k е бројот на зафатени ќелии. Нема дрво, нема сортирање, нема прилагодување — добро прилагодени за множества што брзо се деформираат каде што секој примитив се движи на секој чекор.

\textbf{Опуштени ок-дрва и k-d дрва (octrees and k-d trees)}. Хиерархиски поделби на просторот; корисни кога сцената има многу нерамномерна густина. Изградбата и пречекорувањето се посложени и од BVH и од хеш; ретко се привлечни за работните оптоварувања во оваа теза.




